% Source problem: public/notes/物工/問題/2017ap_exam.pdf, page 3.
\documentclass[11pt,a4paper]{article}
\usepackage[margin=24mm]{geometry}
\usepackage{fontspec}
\usepackage{xeCJK}
\setCJKmainfont{HaranoAjiMincho-Regular.otf}
\usepackage{amsmath,amssymb}
\setlength{\parindent}{1em}
\setlength{\parskip}{0.35em}
\title{物理工学専攻 院試解答集\\2017年 物理学I 第2問}
\author{政岡凜太郎}
\date{}
\begin{document}
\maketitle

\section*{[1.1] 電場と電位}
ガウスの法則より、半径$a$より内側では包む電荷がなく、球殻間だけで電場がある。外側では全電荷がゼロである。電位は無限遠をゼロとして積分する。
\begin{equation}
 E(r)=\begin{cases}0&(r<a),\\ Q_0/(4\pi\varepsilon_0r^2)&(a<r<b),\\0&(r>b),\end{cases}\quad
 \phi(r)=\frac{Q_0}{4\pi\varepsilon_0}\begin{cases}1/a-1/b&(r<a),\\1/r-1/b&(a\le r\le b),\\0&(r>b).\end{cases}
\end{equation}

\section*{[1.2] 静電容量}
球殻間の電位差は$\Delta\phi=Q_0(1/a-1/b)/(4\pi\varepsilon_0)$であるから
\begin{equation}
 C=\frac{Q_0}{\Delta\phi}=\boxed{\frac{4\pi\varepsilon_0ab}{b-a}}.
\end{equation}

\section*{[1.3] 誘電体で満たした場合のエネルギー}
誘電率$\varepsilon$の媒質で満たすと$C_\varepsilon=4\pi\varepsilon ab/(b-a)$である。電荷を$Q_0$に保つので
\begin{equation}
 U=\frac{Q_0^2}{2C_\varepsilon}
 =\boxed{\frac{Q_0^2}{8\pi\varepsilon}\left(\frac1a-\frac1b\right)}.
\end{equation}

\section*{[1.4] 薄い誘電層による容量変化}
真空層$a<r<b-d$と誘電層$b-d<r<b$は直列である。
\begin{align}
 \frac1C&=\frac1{4\pi\varepsilon_0}\left(\frac1a-\frac1{b-d}\right)
 +\frac1{4\pi\varepsilon}\left(\frac1{b-d}-\frac1b\right)\\
 &=\frac1{4\pi\varepsilon_0}\left(\frac1a-\frac1b\right)
 -\frac{d}{4\pi b^2}\left(\frac1{\varepsilon_0}-\frac1\varepsilon\right)+O(d^2).
\end{align}
$C_0=4\pi\varepsilon_0ab/(b-a)$を使って逆数を展開すると
\begin{equation}
 \boxed{C=C_0+\frac{4\pi\varepsilon_0(\varepsilon-\varepsilon_0)}{\varepsilon}\frac{a^2}{(b-a)^2}d+O(d^2)}.
\end{equation}

\section*{[2.1] 定常電流とジュール熱}
媒質内の電場は$E_r=Q_0/(4\pi\varepsilon r^2)$なので、半径によらず電流は
\begin{equation}
 I=4\pi r^2\sigma E_r=\frac{\sigma Q_0}{\varepsilon}>0.
\end{equation}
電位差を電流で割れば抵抗が得られる。
\begin{equation}
 \boxed{\mathcal R=\frac{1}{4\pi\sigma}\left(\frac1a-\frac1b\right)},\qquad
 \boxed{\dot Q_{\rm heat}=I^2\mathcal R=\frac{\sigma Q_0^2}{4\pi\varepsilon^2}\left(\frac1a-\frac1b\right)}.
\end{equation}

\section*{[2.2] 電源を外した後}
時刻$t$の内側球殻の電荷を$Q(t)$とする。[2.1]の電流式に現在の電荷を入れると$I(t)=\sigma Q(t)/\varepsilon$となる。外向きの電流は内側球殻の電荷を減らすので、電荷保存則から$dQ/dt=-I(t)=-\sigma Q(t)/\varepsilon$である。初期条件$Q(0)=Q_0$を使うと
\begin{equation}
 Q(t)=Q_0e^{-\sigma t/\varepsilon}.
\end{equation}
瞬時電力は[2.1]の$Q_0$を$Q(t)$に置き換えたものなので
\begin{equation}
 W(t)=\int_0^t\frac{\sigma Q(t')^2}{4\pi\varepsilon^2}\left(\frac1a-\frac1b\right)dt'
 =\boxed{\frac{Q_0^2}{8\pi\varepsilon}\left(\frac1a-\frac1b\right)\left(1-e^{-2\sigma t/\varepsilon}\right)}.
\end{equation}

\section*{[2.3] 静電エネルギーとの関係}
$t\to\infty$では残留電荷がゼロとなり、発生熱は$W(\infty)=Q_0^2(1/a-1/b)/(8\pi\varepsilon)$となる。これは[1.3]の初期静電エネルギーに等しい。ここでは電源が外れており、媒質の熱以外にエネルギーが移らないためである。

\section*{[3] 二媒質境界の面電荷}
定常状態では各領域の電流密度が$I/(4\pi r^2)$である。従って境界の両側の電場は$E_i=I/(4\pi\sigma_i r_0^2)$。境界面法線を外向きに取ると、自由面電荷密度は法線方向の$\boldsymbol D$の跳びで
\begin{equation}
 \boxed{\sigma_{\rm s}=\varepsilon_2E_2-\varepsilon_1E_1
 =\frac{I}{4\pi r_0^2}\left(\frac{\varepsilon_2}{\sigma_2}-\frac{\varepsilon_1}{\sigma_1}\right)}.
\end{equation}
負値の場合は、境界に負の自由電荷が蓄積している。
\end{document}
